\documentclass[12pt]{article}
\usepackage{amsmath, amsthm, amssymb}
\usepackage[utf8]{inputenc}
\usepackage[margin=1in]{geometry}

\newtheorem{theorem}{Theorem}

\begin{document}

\section{Supremum of a Union of Points and Intervals}
\begin{theorem}
Let $\mathcal{M} = (M, <, \dots)$ be an o-minimal structure.
Let $S \subseteq M$ be a non-empty set that can be expressed as the union of a finite number of points and open intervals.
Then $S$ has a supremum in $M \cup \{+\infty\}$.
\end{theorem}
\begin{proof}
Let $S$ be expressed as $S = P \cup I_1 \cup \dots \cup I_k$, where $P$ is a finite set of points and each $I_j$ is an open interval in $M$.
Because we consider the extended domain $M \cup \{+\infty\}$, every individual component of $S$ has a supremum.
The finite set of points $P$, if non-empty, has a maximum which serves as its supremum in $M$.
For each open interval $I_j$, it can take the form $(a_j, b_j)$, $(-\infty, b_j)$, $(a_j, +\infty)$, or $(-\infty, +\infty)$ for $a_j, b_j \in M$.
If the interval is bounded above, the supremum is precisely its right endpoint, $b_j \in M$. If the interval is unbounded above, its supremum is $+\infty$.
Therefore, $S$ is a finite union of sets, each of which has a well-defined supremum in $M \cup \{+\infty\}$.
The supremum of a finite union of sets, each possessing a supremum, is the maximum of their individual supremums.
Thus, whether $S$ is bounded above or not, it has a supremum in $M \cup \{+\infty\}$.
\end{proof}

\end{document}
